\SourceLeaf{157}
\Needspace{6\baselineskip}\pdfbookmark[3]{II. Bewertung der Verpflichtungen gegenüber gegenwärtigen Teilnehmern}{witwen-II-2-II}
\par\medskip{\raggedright\noindent II.\quad Bewertung des jetzigen Geldwerts der Verpflichtungen der Witwenkasse gegenüber ihren gegenwärtigen Teilnehmern.\par}\smallskip
Wie die Witwen habe ich auch alle Teilnehmer der Witwenkasse in der Reihenfolge ihres Beitritts fortlaufend nummeriert. Ich folge ganz dem an der angeführten Stelle vorgezeichneten Weg und habe zuerst die 42 verheirateten Teilnehmer zu betrachten. Zunächst ist, was gleichsam den Kern der Untersuchung bildet, nach den Grundsätzen der Wahrscheinlichkeitsrechnung und den Sterblichkeitstafeln der gegenwärtige Geldwert sowohl der Beiträge zu ermitteln, zu denen sie verpflichtet sind, als auch der Pensionen, die ihre jetzigen Ehefrauen erhalten werden, falls sie ihre Männer überleben. Ich fasse hier auf einer Seite die Ergebnisse einer langen Arbeit in einer Tabelle zusammen.

\begin{center}
\scriptsize\setlength{\tabcolsep}{1pt}
\begin{tabular}{r|c|r|r|r|r||r|c|r|r|r|r}
&\shortstack{Jetzige\\verhei-\\ratete\\Teil-\\nehmer}&\multicolumn{2}{c|}{Zinsfuß \(3\tfrac{1}{2}\) Prozent}&\multicolumn{2}{c||}{Zinsfuß 4 Prozent}&&\shortstack{Jetzige\\verhei-\\ratete\\Teil-\\nehmer}&\multicolumn{2}{c|}{Zinsfuß \(3\tfrac{1}{2}\) Prozent}&\multicolumn{2}{c}{Zinsfuß 4 Prozent}\\\cline{3-6}\cline{9-12}
Nr.&&\multicolumn{2}{c|}{\shortstack{Jetziger Geldwert\\der}}&\multicolumn{2}{c||}{\shortstack{Jetziger Geldwert\\der}}&Nr.&&\multicolumn{2}{c|}{\shortstack{Jetziger Geldwert\\der}}&\multicolumn{2}{c}{\shortstack{Jetziger Geldwert\\der}}\\
&&Pensionen&Beiträge&Pensionen&Beiträge&&&Pensionen&Beiträge&Pensionen&Beiträge\\\hline
&&Thl.&Thl.&Thl.&Thl.&&&Thl.&Thl.&Thl.&Thl.\\
109&L.&1752&59&1634&58&174&H.&1175&120&1053&114\\
135&O.&1509&85&1386&83&175&R.&1154&142&1007&134\\
138&C.&1846&70&1706&68&177&W.&1117&131&991&124\\
139&U.&937&109&854&104&180&T.&1064&132&944&125\\
140&H.&1451&106&1310&101&183&V.&943&148&824&139\\
141&L.&1091&93&1004&90&184&B.&1023&146&893&137\\
147&R.&1441&105&1301&101&187&W.&917&134&817&127\\
149&G.&1717&101&1545&97&188&S.&867&134&774&127\\
151&O.&1051&130&933&123&189&H.&1138&129&1011&122\\
153&B.&865&109&790&104&190&H.&936&142&825&134\\
155&K.&925&114&840&109&191&D.&892&123&804&117\\
156&S.&853&114&776&109&192&R.&943&139&848&131\\
160&H.&1376&123&1221&117&193&G.&1047&148&907&140\\
161&B.&876&120&795&114&195&D.&981&140&866&132\\
164&W.&1359&122&1209&117&197&R.&916&154&791&145\\
167&S.&995&140&878&132&198&W.&1072&144&936&136\\
168&K.&1345&135&1173&128&199&F.&1127&127&1005&120\\
169&Z.&1118&134&988&127&200&L.&861&153&749&144\\
170&R.&1541&99&1397&96&201&W.&1073&145&935&137\\
171&C.&1193&105&1084&101&203&B.&932&127&835&121\\
173&F.&1120&127&998&121&204&E.&787&146&729&138\\\cline{6-12}
\multicolumn{6}{r}{Summe}&\multicolumn{2}{c|}{Thaler}&47324&5203&42366&4945
\end{tabular}
\end{center}
\noindent\hspace*{\parindent}Den Gesamtwert der Witwenpensionen vergrößere ich wegen der Waisenpensionen um seinen sechsten Teil, nicht so sehr deshalb, weil \textsc{Kritter} in seinem Gutachten dieses Verhältnis angenommen hat, sondern weil es sich sehr nahe aus meiner eigenen Erörterung der Erfahrungen bei unserer Witwenkasse ergibt [S. 154]. Danach wird der Gesamtwert der Beiträge abgezogen. Daraus ergibt sich der reine Wert der Verpflichtung der Kasse gegenüber den 42 verheirateten Mitgliedern. Für die 9 jetzt Unverheirateten wird schließlich ohne Weiteres pro rata ein Betrag hinzugerechnet. Diese Rechnungen sehen dann folgendermaßen aus:
\begin{center}
\footnotesize\setlength{\tabcolsep}{3pt}
\begin{tabular}{l r|r}
&Zinsfuß \(3\tfrac{1}{2}\) Prozent&Zinsfuß 4 Prozent\\
\shortstack[l]{Witwenpensionen der 42\\verheirateten Mitglieder}&47324 Thl.&42366 Thl.\\
Waisenpensionen . . . . . . . . . . . . . . . . . . . . . .&7887 ,, &7061 ,, \\[5pt]
Pensionen . . . . . . . . . . . . . . . . . . . . . . . . . . .&55211 ,, &49427 ,, \\
Beiträge . . . . . . . . . . . . . . . . . . . . . . . . . . . .&5203 ,, &4945
\end{tabular}
\end{center}
\SourceLeaf{158}
\begin{center}
\footnotesize\setlength{\tabcolsep}{3pt}
\begin{tabular}{l r|r}
&Zinsfuß \(3\tfrac{1}{2}\) Prozent&Zinsfuß 4 Prozent\\
\shortstack[l]{Reine Verpflichtung der Kasse gegen\\42 verheiratete Mitglieder}&50008 Thl.&44482 Thl.\\
\shortstack[l]{Danach verhältnismäßig gegen\\9 unverheiratete Mitglieder . . . .}&10716 ,, &9532 ,, \\[5pt]
\shortstack[l]{Gesamtwert der zweiten Rubrik\\ . . . . . . . . . . . . . . . . . . . . . . . .}&60724 ,, &54014 ,,
\end{tabular}
\end{center}
\noindent\hspace*{\parindent}Dass einer solchen Bewertung nicht ganz dieselbe Zuverlässigkeit zuerkannt werden kann wie der Berechnung der Verpflichtungen unter der ersten Rubrik, habe ich bereits in meinem mehrfach erwähnten Gutachten vorausgeschickt. Bei unserer Witwenkasse liegen nämlich mehrere Umstände vor, deren Wirkung überhaupt nicht im Voraus berechnet werden kann. Weil für die frühere Zeit eine Buchführung fehlt, wie sie S. [145, unten] erwähnt ist, lässt sich diese Wirkung nicht einmal schätzen. Als den einflussreichsten dieser Umstände bezeichne ich die Wiederverheiratung von Mitgliedern nach dem Tod ihrer Ehefrauen. Dass die Frauen vor den Männern sterben, gehört nämlich bereits zu den in der Rechnung berücksichtigten Möglichkeiten. Jede Wiederverheiratung ist daher eine außerhalb der Rechnung liegende, neu hinzukommende Belastung des Kassenkontos. Auf der anderen Seite kommt diesem Konto jeder sonstige Abgang eines verheirateten Mitglieds, etwa durch eine Berufung an einen anderen Ort, als eine nicht veranschlagte Entlastung zugute. In dieselbe Kategorie gehört auch, dass solche verwitweten Mitglieder, die nicht wieder heiraten, nach dem Tod ihrer Frauen doch wenigstens eine Zeit lang ihre Beiträge weiterzuzahlen pflegen; dies ist allerdings von geringer Bedeutung. Schließlich ist die Rechnung für die derzeit unverheirateten Teilnehmer nach dem Durchschnittsergebnis für die Verheirateten vielleicht etwas zu hoch. Der jeweilige Bestand der unverheirateten Mitglieder wird gewöhnlich so zusammengesetzt sein, dass man mit Wahrscheinlichkeit annehmen kann, der eine oder andere werde überhaupt nicht heiraten. Dies wird jedoch zumindest teilweise wieder dadurch aufgewogen, dass man, je später die Heirat erfolgt, umso eher damit rechnen muss, dass die Frau ihren Mann überlebt und lange im Witwenstand bleibt.

Obwohl es unmöglich ist, die Wirkung dieser verschiedenen Umstände im Voraus in Zahlen zu veranschlagen, darf man doch für wahrscheinlich halten, dass sie sich größtenteils gegenseitig ausgleichen.

